\section{Introduction}

Rare earths and strategic metals connect commodity pricing with industrial security. Magnets, batteries, electricity networks and manufacturing equipment depend on access to particular materials and processing stages. Geopolitical tensions can therefore change expected costs and availability. The important empirical question is whether a systematic measure of those tensions explains prices and helps an investor make better decisions.

The supply-chain story does not give a unique investment prediction. A possible shortage can raise the price of a material while increasing the costs faced by firms that use it. A producer may benefit from a higher selling price, but face financing problems or new restrictions of its own. A fund holding companies across several countries and stages of production combines these exposures. The price of a magnet material and the return on a producer-equity portfolio can consequently respond differently to the same news.

The paper asks three questions. First, does a change in geopolitical risk have a different price association when risk was already elevated? Second, does that association remain after accounting for economic activity, the dollar and other commodity prices, and how does it evolve over subsequent months? Third, do geopolitical signals add investment value relative to standard price-based rules and market risk factors? The distinction between explaining a current price change and predicting a future attainable return is central to all three questions.

The main physical series is praseodymium–neodymium mixed metal, abbreviated PrNd. Monthly quote observations cover January 2010–January 2026. The selected interaction model contains 191 returns after allowing for lagged information. Other rare earth quotes, eight IMF material benchmarks and sector equity funds provide additional context. The geopolitical measure is the Caldara–Iacoviello GPR index, particularly its China-related series \citep{caldara2022}.

Investment evaluation uses REMX, LIT and COPX: funds associated with rare earth and strategic materials, the lithium and battery value chain, and copper mining. A chronological exercise starts in March 2018, after the original working paper, with initial estimation ending in December 2017. It uses a consistent revised risk history and is labelled accordingly. A separate 2022–2026 reconstruction reads risk values and refits forecasts from the actual archived release used at each decision. Earlier observations remain useful for estimating historical relationships; their availability as a historical series is distinct from the availability of a live investment signal.

The price results survive a meaningful extension. The selected high-state PrNd slope changes from 3.004 to 3.001 when lagged measures of Chinese activity, global activity, the broad dollar and other commodity returns are included. The evidence for a difference between the two risk states is less precise. Local projections give the clearest positive association in the current month and much wider intervals for later cumulative changes. This pattern is consistent with a contemporaneous price response that need not leave a simple next-month profit opportunity.

The investment results are heterogeneous. Over the revised-data exercise from 2018, avoiding high and rising China-related risk earns \OosLitAvoidCagr\% annually in LIT and \OosCopxAvoidCagr\% in COPX after assumed costs. Standard trend rules also earn positive returns, and conclusions change across funds, execution conventions and calibration windows. Six-factor attribution and nested forecast comparisons do not establish a reliable incremental geopolitical premium. Archived risk data further show how two revised classifications remove losing REMX exposures from a later historical reconstruction.

The central proposition is therefore that \textit{the pricing and investment implications of geopolitical risk are state-dependent, commodity-specific, and sensitive to the information available at the time of the decision}. The contribution is to examine these dimensions together. It combines a transparent conditional price model, economic controls and dynamic responses with benchmarked investment tests and a decision-level account of data revisions. The analysis is exploratory, with its complete specification families and unsuccessful rules retained for replication.

\section{Economic mechanisms and related literature}

A supply channel predicts higher prices when an escalation threatens availability and substitution is difficult. Buyers may accelerate orders or increase precautionary stocks before an actual restriction takes effect. A demand channel can work in the opposite direction: uncertainty can delay industrial investment and reduce expected purchases. Exchange rates, financing conditions and other commodity costs can move with geopolitical news. The sign of an average regression coefficient reflects the combined operation of these forces.

The pre-existing risk environment can matter. During a calm period, a warning may have little effect on purchasing plans. After persistent tension, an additional warning may trigger stockpiling or a revision of quoted availability. An interaction between a new risk change and the previous risk level is a simple way to test this idea. It detects a conditional association; inventories, contracts and lost production are not directly observed here, so it does not identify the supply mechanism by itself.

\citet{caldara2022} establish a news-based framework for measuring adverse geopolitical events and threats. \citet{fan2023} relate geopolitical risk to rare earth trade using Japanese import data, finding higher import unit prices and lower import values, with China-related risk particularly relevant. An import unit value can also change with product composition and trading partners. It therefore motivates the supply-chain hypothesis without fixing the response of a particular quote or equity fund.

\citet{li2023} examine rare earth prices, geopolitical risk and economic activity across time and frequencies. \citet{depraiter2025} study connections with clean-energy investment, and \citet{magnets2025} examine Chinese permanent-magnet exports. The industrial context includes concentrated critical-mineral supply chains and trade restrictions \citep{iea2025}, while regulatory definitions distinguish particular strategic products and uses \citep{eu2024}. These studies support attention to the measured product and its place in the value chain rather than a single coefficient for all materials.

The investment comparison draws on a separate literature. Time-series momentum uses an asset's own past return to determine exposure \citep{moskowitz2012}. Factor models distinguish average returns from compensation associated with familiar equity exposures \citep{fama2015}. Return forecasting must also compete with simple alternatives \citep{campbell2008}. These comparisons make the economic question more demanding: GPR should contribute something beyond a sector's market exposure or a standard price trend, rather than merely improve on one weak passive holding.
